% RESEARCH_DRAFT_SHA256: a80532f361bbc525ccdff8294302bcdc15c4d5de2e439816200bde736826051f
% Layout and source-reference wrappers expanded for the web renderer.
\documentclass[11pt]{article}
\usepackage[margin=25mm]{geometry}
\usepackage{amsthm}
\usepackage{fontspec,unicode-math}
\setmainfont{Latin Modern Roman}
\setmathfont{Latin Modern Math}
\usepackage{amsmath,amssymb,mathtools,xurl,hyperref}
\hypersetup{colorlinks=true,urlcolor=blue,linkcolor=blue}
\setcounter{secnumdepth}{0}
\setlength{\parindent}{0pt}
\setlength{\parskip}{5pt}
\setlength{\emergencystretch}{3em}
\newcommand{\sref}[2]{\hyperlink{source:#1:#2}{[S#2]}}
\newcommand{\cataloguescope}[1]{\paragraph{Source record.}#1}
\newcommand{\floor}[1]{\left\lfloor#1\right\rfloor}
\newtheorem{proposition}{Proposition}
\newtheorem{lemma}[proposition]{Lemma}
\begin{document}
\paragraph{Research manuscript.} OpenAI \texttt{gpt-6-astra}, reasoning \texttt{ultra}; 11 October 2026.
\paragraph{Preserved source body.} The complete catalogue mathematical text and sources follow with layout and source links expanded, including source qualifications and its historical preparation notes. New research begins at the section headed Research attempt.

\section{Neural Approximation of Strongly Convex Drifts with Linear Error}\label{statement}

\subsection{Definitions and mathematical statement}

\paragraph{Shared notation.}Standard number systems are $\mathbb N,\mathbb Z,\mathbb Q,\mathbb R,\mathbb C$. Specialized notation and index ranges follow the cited source and the definitions below.


Under the Euclidean convention, a differentiable $V:\mathbb R^d\to\mathbb R$ is $\mu$-strongly convex when $\mu>0$ and $V(y)\ge V(x)+\langle\nabla V(x),y-x\rangle+(\mu/2)\|y-x\|^2$ for all $x,y$. Its gradient is $L$-Lipschitz when $\|\nabla V(x)-\nabla V(y)\|\le L\|x-y\|$. The drift is $b(x)=-\nabla V(x)$; a neural approximation $N$ has pointwise error $\|N(x)-b(x)\|$. A feed-forward network composes affine maps with specified activation functions.
\paragraph{Record scope and qualifications.}
The catalogue title is retained as metadata. Named certificates, contracts, bounds or frameworks in that title are not asserted as results of the original source extract.
Source-record limitation: For an $L$-Lipschitz gradient, the zero approximation has the affine-linear error bound $\|\nabla V(x)\|\le\|\nabla V(0)\|+L\|x\|$. The constant network $N(x)=-\nabla V(0)$ even satisfies $\|N(x)+\nabla V(x)\|\le L\|x\|$. Thus merely requiring linear growth does not specify a nontrivial approximation target. The cleaned question omits the accuracy, norm, and coefficient constraints that would distinguish the intended result.
\paragraph{Statement recorded in the supplied source.}
If $V$ is strongly convex with Lipschitz gradient, does there exist a neural network approximating $-\nabla V$ with error growing only linearly? \href{https://doi.org/10.4171/owr/2023/48}{[S2]}
\paragraph{Source status.}
The supplied catalogue labels this source record partially solved; that metadata was checked against the supplied snapshot on 2026-10-11. The label is preserved for provenance and is not an independent verification of present open status or mathematical correctness. The catalogue literature assessment is dated 2026-08-22.
\subsection{Short English statement}
Approximate the negative gradient of a strongly convex potential by a neural network with a prescribed bound on how the error grows with distance.
\subsection{Sources}
\begin{itemize}
\item[S1] Ulam AI and the UnsolvedMath Contributors, UnsolvedMath: A Curated Collection of Open Mathematics Problems, catalogue entry. Catalogue curation and metadata: CC BY 4.0. The source extract is retained; definitions, English summaries and qualifications are editorial additions.
\url{https://huggingface.co/datasets/ulamai/UnsolvedMath}.
\item[S2] Mini-Workshop: Nonlinear Approximation of High-dimensional Functions in Scientific Computing (2024), Oberwolfach Reports, DOI 10.4171/owr/2023/48, pp. 2771--2808
\url{https://doi.org/10.4171/owr/2023/48}.
\item[S3] Eigel, Miranda, Schuette, and Sommer, WIAS Preprint 3077
\url{https://www.wias-berlin.de/preprint/3077/wias_preprints_3077.pdf}.
\item Source-specific Oberwolfach Reports licences, including CC BY-SA 4.0 where applicable, are retained; see THIRD\_PARTY\_NOTICES.md.
\end{itemize}
\paragraph{Preparation and review.}
The original source and UnsolvedMath contributors retain attribution for the source extract. Codex (GPT-6; exact serving revision not exposed) assisted with definitions, English summaries, scope notes and formatting. No new human mathematical authorship or independently verified solution is asserted. Catalogue identity checked against the complete supplied registry and current repository. Distinct upstream records, including repeated questions and broad or damaged extracts, are retained for complete coverage of the source CS category. A repeated record is not a new independent conjecture. Complete mathematical-equivalence and current-literature review has not been performed. Preparation and final review follow contribution prompts C and D, with source gaps disclosed. Structural and rendering checks do not establish a new theorem or open status.
\label{end:statement}

\section{Research attempt}
\subsection{Literature, source context and scope}
The full original catalogue body and source qualifications are preserved above. The relevant primary workshop passages were read in the cached publisher report. Historical claims and subsequent attributed results are distinguished from the mathematical work below. A restricted example, standard lemma or known theorem is not counted as new progress on the entire original target. Missing assumptions and repeated or rhetorical extracts are not silently turned into different solved problems.

\subsection{The error coefficient and activation matter}
As the preserved definition notes, a constant network already has an error bounded by $L\|x\|$ for an $L$-Lipschitz gradient. The primary WIAS paper makes a stronger assumption in its Assumption1.2(i): for each requested $\varepsilon>0$, global error is at most $\varepsilon(1+\|x\|)$. Its local-approximation route with the extra condition $M<\sqrt2m$ is a distinct sampling theorem, not a claim that arbitrary potentials satisfy that stronger assumption. We test the stronger global statement for finite ReLU networks, an explicitly specified activation class.
Define in one dimension
\[
g(x)=2x+\frac{x}{8}\sin\!\left(\tfrac12\log(1+x^2)\right),
\qquad V(x)=\int_0^xg(t)\,dt.
\]
\begin{proposition}
The potential $V$ is $7/4$-strongly convex and has $9/4$-Lipschitz gradient. Nevertheless no finite feed-forward ReLU network $N:\mathbb R\to\mathbb R$ satisfies
\[
|N(x)+V'(x)|\le\varepsilon(1+|x|)\quad\hbox{for all }x\in\mathbb R
\]
when $0<\varepsilon<1/8$.
\end{proposition}
\begin{proof}
Put $\theta(x)=\tfrac12\log(1+x^2)$. Differentiation gives
\[
g'(x)=2+\tfrac18\sin\theta(x)+\tfrac18\frac{x^2}{1+x^2}\cos\theta(x),
\]
which lies in $[7/4,9/4]$. The mean value theorem gives the gradient Lipschitz bound. Integrating $V''\ge7/4$ twice along the segment from $x$ to $y$ yields $V(y)\ge V(x)+V'(x)(y-x)+(7/8)(y-x)^2$, including either ordering of $x,y$.
Every finite one-input ReLU network is continuous piecewise affine with finitely many breakpoints. This follows by induction over its layers: finite affine combinations preserve a common finite interval partition, and on each interval ReLU of an affine function introduces at most one new breakpoint unless that function is identically zero. Hence on some positive tail $N(x)=ax+b$.
For each large integer $j$, choose $x_j>0$ with $\theta(x_j)=\pi/2+2\pi j$ and $z_j>0$ with $\theta(z_j)=3\pi/2+2\pi j$; explicitly these are $\sqrt{e^{2\theta}-1}$ and tend to infinity. Then $g(x_j)/x_j=2+1/8$ and $g(z_j)/z_j=2-1/8$. Dividing the assumed error bound by the input and taking these two limits gives
\[
|a+2+1/8|\le\varepsilon,\qquad |a+2-1/8|\le\varepsilon.
\]
The triangle inequality would imply $1/4\le2\varepsilon$, a contradiction.
\end{proof}
The constants even satisfy $(9/4)/(7/4)=9/7<\sqrt2$, emphasizing why the WIAS sampling theorem and the arbitrary-small-global-slope assumption must not be conflated. This is not a contradiction to its theorem. Other activations may have different tails; the argument does not cover them. The unrestricted wording ``linear error growth'' remains trivially satisfied with a sufficiently large coefficient, so this counterexample does not refute that raw statement.
\subsection{Verification and remaining gaps}
The derivative formula and bounds were checked on numerical grids, the two limiting slopes were checked at moderate exact phase choices, and the interval contradiction was audited symbolically. The tail proof is exact and does not depend on numerical extrapolation. The result blocks one natural strengthened ReLU formulation, but the original record does not fix that formulation or activation, approximation complexity, or accuracy coefficient. No full-source solution or new global sampling theorem is claimed.

\section{References and source attribution}
Primary workshop report: \url{https://doi.org/10.4171/owr/2023/48}. Relevant passages read: 2023/48 pp. 2792--2793, Langevin neural approximation, L2 law-dependent error and linear-growth conjecture.

\url{https://www.wias-berlin.de/preprint/3077/wias_preprints_3077.pdf}. Read primary Assumptions1.1--1.2, Theorem1.3 and explanatory discussion on PDF pp.3--6; arbitrary epsilon global assumption and separate condition-number sampling route are distinguished.

All retrievals were checked on 11 October 2026. Original question attribution remains in the preserved Sources section. External theorems are attributed and not claimed as new discoveries or fully verified merely by citation.


\paragraph{Integration review.} Codex (GPT-6; exact serving revision not exposed) reviewed the written arguments and expanded layout and source-reference wrappers. Repository metadata and the production web renderer were checked using Python 3.13 and JavaScript with MathJax 3.2.2. This model-assisted review is not independent human proof verification; no PDF compilation or formal proof checking is asserted.

\section{Completion Estimate}
\noindent\textbf{COMPLETION: 0\%}
No substantiated new progress toward the entire original target is claimed. The worked results and failed route are delimited above; a known result, illustrative model, or auxiliary obstruction is not counted as original-target completion.

\section{Final status and attribution}
\textbf{Attempt status: unresolved.} This describes this research attempt, not a new ruling on the source catalogue's status or on external literature claims.
The mathematics and this self-audit were generated by OpenAI \texttt{gpt-6-astra}, reasoning \texttt{ultra}, on 11 October 2026. Codex (GPT-6; exact serving revision not exposed) orchestrated the batch, read and reviewed its written mathematical arguments, and contributed formatting/repository review. No human mathematical authorship, independent human proof checking, or formal verification is claimed. The historical preparation paragraphs in the reproduced catalogue body concern that body only.

Contribution prompt A was applied to the complete definition and supplied source record, followed by the adversarial checks of prompt B. Local reproduction files are \texttt{batch025/checks.py}, \texttt{computations.json}, and the per-ID audit JSON. They were not public links at preparation time. Numerical checks are finite and do not prove unrestricted results. No LaTeX engine was found locally; compilation was not performed by this research agent.

\end{document}
